\section{Open challenges and research directions}
\label{sec:challenges}

The synthesis in Sections~\ref{sec:integration}--\ref{sec:applications}
identifies recurring gaps that cut across surrogate families,
integration patterns and application domains. We collect them here as
six research directions. The intent is not to cover every issue raised
in the literature but to pick the gaps that are repeatedly reported
\emph{and} that would change the way surrogate-assisted MES studies
are designed and reported.

\subsection{Decision-aware benchmarking}
\label{sec:ch:bench}

The most consistent gap is the absence of decision-aware benchmarking.
Studies report regression metrics on a held-out set and infer
decision quality only indirectly; speed-ups are reported routinely,
the corresponding regret distribution rarely
\cite{Mylonopoulos202332697,Ruan2021221,Khaloie2025,Lim2025,Starke2025214,Perera2019191}.
A decision-aware benchmark would fix the host optimization problem
(security-constrained dispatch, AC OPF, multi-objective MES design)
together with a public set of training and stress-test instances and
require all submitted surrogate methods to report regret, feasibility
rate, MIP-gap distribution and worst-case violation under a common
solver-time budget. Recent learning-to-optimize work for security-
constrained dispatch and unit commitment establishes the necessary
building blocks but stops short of a shared benchmark
\cite{Wu2022231,Chen20244723,Chen2022,Wu2024391,Liang20246508}. The
construction of such a benchmark is plausibly the highest-leverage
methodological intervention because it would make subsequent
contributions directly comparable.

\subsection{Calibrated uncertainty for surrogate-assisted MOO}
\label{sec:ch:uq}

Surrogate-assisted multi-objective optimization rests on calibrated
posterior uncertainty: the acquisition function in P2 instances and
the chance constraint in P5 instances both depend on it. Empirical
evidence in the screened pool indicates that posterior calibration is
checked far less often than mean-square error, and that calibration
breakdown is the typical failure mode in operating regions that the
training set did not cover
\cite{Hu2020,Hu2021671,Volodina2022,Lawson201647,Hu2024,Buchanan2024,Subedi2026157}.
Recent advances in deep ensembles, Bayesian last layers, conformal
prediction and chance-constrained polynomial-chaos OPF provide
practical instruments
\cite{Chen2023657,Liang20246508,Muhlpfordt20194806,Koirala20242284,Xu20212696}.
The open question is how to combine these instruments with the
acquisition functions used in surrogate-assisted MOO so that the
expected hypervolume improvement reflects the actual posterior
spread on energy-system inputs, not a nominal Gaussian assumption.

\subsection{Out-of-distribution generalisation across capacity mixes}
\label{sec:ch:ood}

Energy-system planning studies are intrinsically counterfactual:
results matter precisely because they describe future capacity mixes,
markets and weather patterns. The training distribution of an energy-
system surrogate, however, is set by historical operation or by the
optimizer that produced the synthetic data. Out-of-distribution
robustness is therefore a planning-relevance criterion, not just a
machine-learning hygiene requirement
\cite{Park2019,Reynolds2017704,Subedi2026157,Khaloie2025,Lim2025,Starke2025214,Mylonopoulos202332697}.
Two complementary research lines stand out. The first is invariant or
physics-informed feature design that ties the surrogate to
conservation laws (power balance, energy balance, mass flow) so that
the surrogate respects them on inputs it has not seen
\cite{Park2019,Su2022315,Liu20223726,Petrova2025,Du2025}. The second
is transfer-learning protocols that explicitly fine-tune a base
surrogate to a new climate, year or capacity mix using a small
high-fidelity sample
\cite{Chen202514792,De Castro20247710,Khayambashi202453}. Both lines
need a shared evaluation protocol that distinguishes within-
distribution from out-of-distribution performance.

\subsection{Standardised public datasets and reproducibility}
\label{sec:ch:data}

The reproducibility audit in Section~\ref{sec:val:repro} indicates
that standardised public datasets are scarce, particularly for the
multi-energy and district-heating settings that are central to
this review
\cite{Khaloie2025,Lim2025,Starke2025214,Ruan2021221,Mylonopoulos202332697,Perera2019191,Schulte2016104,Chaudhry2024,Bianchi2020244}.
Calibration-focused communities that depend on external reference
data report reproducible artefacts more consistently
\cite{Azevedo2021630,Khazeiynasab2021,Zhang2025__7,Ullah20255970}, but
the bulk of MES surrogate studies still ship one-off case studies. A
public dataset that pairs a high-fidelity district-heating, building
or process-energy simulator with a published optimization problem,
the corresponding training set, the train/validation/test split and
solver settings would change the cost--benefit balance of releasing
reproducible studies. Scenario-rich datasets that already exist for
load and renewable forecasting
\cite{Eseye201991463,Faraji2020157284,Hanifi2024,Sideratos2020,Tian2020,Madhukumar20228891}
provide partial templates for what such a release should contain.

\subsection{MCDM under surrogate uncertainty}
\label{sec:ch:mcdm}

The literature on multi-criteria decision making after surrogate-
assisted MOO is sparse despite the central role that posterior
decision making plays in the design of multi-energy systems. TOPSIS,
VIKOR, AHP and PROMETHEE have been applied to surrogate-derived Pareto
fronts in a small number of studies, typically without an explicit
treatment of how surrogate uncertainty propagates through the MCDM
ranking \cite{Roy2020,Roy2023}. Recent multi-objective design work
selects compromise solutions on Pareto fronts produced by NSGA-style
algorithms but treats the front as a deterministic object
\cite{Hussien202415075,Liang2024__3,Hussien202190577,Dong2024,Salahinezhad2025,Cao2017,Zang2023813,Reis2025,Katkar2025,Guan2025,bian_multi-energy_2025,chen_multi-objective_2025-1,cortes-caicedo_multi-objective_2025,tezer_multi-objective_2025,tian_building_2025}.
A consistent treatment would propagate the posterior uncertainty
through the MCDM step and report a distribution over rankings rather
than a single ranking; the lack of such a protocol is, in our view,
the principal methodological gap on the decision-making side of the
surrogate $\times$ MOO $\times$ MES intersection.

\subsection{Open-source libraries and toolchain integration}
\label{sec:ch:tooling}

The screened pool contains several mature surrogate-and-optimisation
toolchains -- polynomial-chaos OPF, learning-to-optimize for unit
commitment, surrogate-assisted evolutionary frameworks for hybrid
renewable systems -- but each toolchain typically reimplements its
own data pipeline, validation harness and reporting interface
\cite{Muhlpfordt20194806,Mitrovic2023__2,Mitrovic2023__3,Wu2022231,Chen20244723,Chen2022,Roth2019214,Hirvonen2017323,Schulte2020}.
A shared open-source library that provides (i) a registry of host
optimization problems, (ii) a uniform interface for surrogate
training and prediction, (iii) the validation harness implied by
Section~\ref{sec:validation}, and (iv) a reporting layer with
decision-aware metrics would lower the cost of contributing new
surrogate variants and would make the resulting evidence directly
comparable. Existing energy-system optimisation frameworks (PyPSA,
Calliope, oemof, SpineOpt, Backbone) provide the host side; the
surrogate side is still fragmented.
